Una olla de sopa se enfría en aire a $20\,{}^\circ\mathrm C$. La pérdida de calor por unidad de tiempo es proporcional a la diferencia entre la temperatura de la sopa y la del aire. Suponga temperatura uniforme, capacidad calorífica constante y ausencia de aporte de calor. Para el cálculo, tome como temperatura inicial de ebullición $100\,{}^\circ\mathrm C$. Al cabo de $30$ minutos se mide $35\,{}^\circ\mathrm C$.

\begin{enumerate}[label=(\alph*),ref=\alph*,leftmargin=*]
\item Formule la ecuación e indique el signo de la constante.
\item Determine la temperatura a los $60$ minutos y el descenso durante los segundos $30$ minutos.
\end{enumerate}
